\section{Proof of Lemma~\ref{lem-geometric-control} and supporting estimates for Theorem~\ref{thm-stability}}\label{sec:9}

We now supply the proof of Lemma~\ref{lem-geometric-control} deferred in Section~\ref{sec:6.1}, and the
auxiliary shape and center estimates deferred in the proof of Theorem~\ref{thm-stability}.
Sections~\ref{sec:9.2}--\ref{sec:9.3} prove the global and local parts of Lemma~\ref{lem-geometric-control}, respectively;
Section~\ref{sec:9.4} supplies the remaining inputs to Theorem~\ref{thm-stability}. The quantities have
three distinct roles: $\epsilon$ measures the total logarithmic volume deficit,
$\chi$ measures departure of the contacts from the equality set, and
$\sigma=2\chi/[d(2+b)]$ is the same contact defect in the block normalization.
We first turn $\chi$ into matrix and contact-distance bounds. We then improve
the determinant lower bound near equality and transfer the resulting matrix
control to the ellipsoid's shape and center.

The constants record successive sufficient estimates in this chain.
The threshold in Lemma~\ref{lem-local} ensures that a quadratic logarithm bound applies;
Section~\ref{sec:9.4} uses fixed spectral bounds to control inversion. The final
coefficients $30$ and $50$ collect those estimates and are not asserted to
be optimal stability constants.

\subsection{Complementarity and matrix matching}\label{sec:9.1}

Use the primal and dual quantities from \eqref{eq:4.1}--\eqref{eq:4.2}. The following identities
also give the weighted multiplier estimate mentioned in Section~\ref{sec:4.3}.
Here $(\tau,x)$ denotes a generic scalar--vector cone pair; the indexed
$(\tau_r,x_r)$ are the support slack and vector from \eqref{eq:4.1}.

\begin{proposition}[{quantitative complementarity and matrix matching}]\label{prop-gap-control}
For cone-feasible $(\alpha,q)$ and $(\tau,x)$,
\begin{equation}\label{eq:9.1}
2\alpha\tau(\alpha\tau-q^Tx)
=\|\alpha x-\tau q\|^2
+\alpha^2(\tau^2-\|x\|^2)+\tau^2(\alpha^2-\|q\|^2).
\end{equation}

If $G\le\epsilon\le1/4$, then
\begin{equation}\label{eq:9.2}
\|Y-I\|_F\le2\sqrt\epsilon,\qquad
\sum_{\alpha_r>0}\alpha_r\tau_r
\left\|\frac{x_r}{\tau_r}-\frac{q_r}{\alpha_r}\right\|^2
\le2\epsilon.
\end{equation}

Here zero-normal redundant inequalities are removed. The first estimate also holds with $Y$ replaced by $A^{1/2}MA^{1/2}$.

\end{proposition}

\begin{proof}
The sum-of-squares identity converts a small cone pairing into
closeness of the normalized cone directions. Expanding the squares proves
\eqref{eq:9.1}, and cone feasibility makes every term on its right nonnegative.
For a nonzero normal, $A\succ0$ and primal feasibility imply
$\tau_r\ge\|Ag_r\|>0$. Divide \eqref{eq:9.1} by $\alpha_r\tau_r$ when $\alpha_r>0$
and sum. The total cone gap is at most $\epsilon$, giving the second bound.

The matrix estimate is Lemma~\ref{lem-matching}:
$\|Y-I\|_F\le2\sqrt G\le2\sqrt\epsilon$.
The symmetric positive-definite matrices $M^{1/2}AM^{1/2}$ and
$A^{1/2}MA^{1/2}$ have the same eigenvalues, so their Frobenius distances
from the identity agree. This proves the final assertion. 
\end{proof}

\subsection{Global part of Lemma~\ref{lem-geometric-control}: geometric residuals}\label{sec:9.2}

We return to the block and contact notation of Section~\ref{sec:8}. The aim is to make
the scalar defect $\chi$ useful in two ways: as a bound on the certificate's
matrix error, and as a bound on the distance of the contacts from the equality
set. For the matrix bound, we seek a residual for each contact whose squared
norm is proportional to that contact's summand in $\chi$. Isotropy will then
combine these residuals without a dimension factor.

\begin{lemma}[{geometric residual}]\label{lem-residual}
For the explicit certificate,
\begin{equation}\label{eq:9.3}
\|Q-M_0\|_F^2\le\frac{1-b}{d}\chi,\qquad
\|\widehat F-I\|_F^2\le\frac\chi d,
\qquad
\sum_i\lambda_i\operatorname{dist}(u_i,\mathcal Z_b)^2
\le\frac{32}{3}\chi.
\end{equation}

\end{lemma}

\begin{proof}
\textbf{Step 1: express the certificate error as contact residuals.} Subtract the model certificate from $Q$ and use the zero first moment to add
a term that vanishes in the aggregate. This makes the individual residual
vanish at both the pole and the model latitude.
For $u=sv+y$ with $y\perp v$, the resulting residual is
\begin{equation}\label{eq:9.4}
\xi(u)=-\frac{s-b}{d(1-as)}
\bigl[(1-b)(s+1)v+\sqrt{1-b}\,y\bigr].
\end{equation}

Expansion of $Q-M_0$, followed by the first and second John moments, gives
\begin{equation}\label{eq:9.5}
Q-M_0=\sum_i\lambda_i u_i\xi(u_i)^T,
\qquad
\|\xi(u)\|^2=\frac{1-b}{d}\frac{(1+s)(s-b)^2}{1-as}.
\end{equation}

For clarity, the axial coefficient before simplification is
$f(s)(s-a)-(2-t)s+b$, and the transverse coefficient is
$kf(s)-\sqrt{1-b}$. The added constant vector $bv$ contributes zero to the aggregate by balance. Their simplification gives \eqref{eq:9.4}, and orthogonality gives the norm identity.

\textbf{Step 2: use isotropy to bound the total matrix error.} The matrix with
columns $\sqrt{\lambda_i}u_i$ has orthonormal rows. Its Bessel inequality
bounds the squared Frobenius norm of the aggregate by the sum of the
weighted squared residuals. Applying this to \eqref{eq:9.5} proves the first estimate.
To transfer it to the normalized matrix, use
$Q-M_0=F_0^{1/2}(\widehat F-I)F_0^{1/2}D$.
The product of the least singular values of the two outside factors is
$\sqrt{1-b}$, giving the second estimate. The two matrix constants in \eqref{eq:9.3}
thus come from the exact residual identity and the fixed normalization.

\textbf{Step 3: compare the scalar residual with geometric distance.}
The equality set consists of a pole and a latitude sphere. We choose an
explicit comparison point in one of these two pieces. It suffices to show
\begin{equation}\label{eq:9.6}
\frac{(1+s)(s-b)^2}{1-as}
\ge\frac3{32}\operatorname{dist}(u,\mathcal Z_b)^2.
\end{equation}

The split at $s=-1/2$ gives simple uniform bounds in both regions.
If $s\le-1/2$, use the pole, whose squared distance is $2(1+s)$;
then $(s-b)^2\ge1/4$ and $1-as\le4/3$.
If $s\ge-1/2$, choose the point on the latitude with the same transverse
direction, choosing any unit transverse vector if $y=0$. Its squared distance is
\begin{equation}\label{eq:9.7}
(s-b)^2+\bigl(\sqrt{1-s^2}-\sqrt{1-b^2}\bigr)^2\le4(s-b)^2.
\end{equation}

Indeed, rationalization bounds the squared root difference by
$(s-b)^2(s+b)^2/(1-b^2)\le3(s-b)^2$.
Also $(1+s)/(1-as)\ge3/8$. These two comparisons give the same coefficient
$3/32$ in \eqref{eq:9.6}; its reciprocal is the contact-distance constant in \eqref{eq:9.3}.
Summing with the John weights finishes the proof. 
\end{proof}

\textbf{Global part of Lemma~\ref{lem-geometric-control} proved.} The first and third inequalities in
\eqref{eq:9.3} are exactly the two estimates in \eqref{eq:6.1}, since $d=2-b$.
The local implication in Lemma~\ref{lem-geometric-control} remains to be proved in Section~\ref{sec:9.3}.

\textbf{What the global estimates provide.} Combining Theorem~\ref{thm-determinant}, Proposition~\ref{prop-gap},
and Lemma~\ref{lem-residual} gives
\begin{equation}\label{eq:9.8}
\chi\le1500\epsilon,\qquad
\sum_i\lambda_i\operatorname{dist}(u_i,\mathcal Z_b)^2\le16000\epsilon,
\qquad
\log\det\widehat F\ge\frac1{1000}\|\widehat F-I\|_F^2.
\end{equation}

\subsection{Local part of Lemma~\ref{lem-geometric-control}: coercivity near equality}\label{sec:9.3}

The global estimate is enough to place a nearly extremal ellipsoid in a small
neighborhood of the model. Here we improve its coefficient within that
neighborhood. The idea is that the positive linear term
$\operatorname{tr}(\widehat F-I)=\sigma-x$ should exceed the quadratic remainder
in $\log\det\widehat F$. We need to bound the whole perturbation, including
$j$, tightly enough for this difference to remain positive.

The threshold $10^{-3}$ below is a sufficient size condition for that
comparison. The intermediate bounds $1/40$ and $1/30$ specify the small
spectral intervals used in the proof; the useful conclusion is the final
lower bound $\sigma/16$.

\begin{lemma}[{local coercivity}]\label{lem-local}
If $\sigma\le10^{-3}$, then
\begin{equation}\label{eq:9.9}
\log\det\widehat F\ge\frac\sigma{16}\ge\frac\chi{32}.
\end{equation}

\end{lemma}

\begin{proof}
If $\sigma=0$, the assertion follows from $\widehat F=I$.
Suppose $\sigma>0$.

\textbf{Step 1: confine the transverse eigenvalues.} We first use Lemma~\ref{lem-variational} to
control $E$ and $j$ in the metric determined by $K_E$.
Here $0\le L\le d(1-b)/(d+b^2)<1$ by its definition and
$(d+b^2)-d(1-b)=2b>0$. Lemma~\ref{lem-variational} therefore gives
\begin{equation}\label{eq:9.10}
A_E+\|K_E^{-1}j\|^2\le\frac\sigma2.
\end{equation}

Every eigenvalue $e>-1/2$ of $E$ therefore satisfies
$e^2/(1+2e)\le1/2000$. This scalar inequality implies $|e|\le1/40$:
the quotient increases as $e$ moves away from zero on either side, and its
values at $\pm1/40$ already exceed $1/2000$.
The choice $1/40$ gives the simple interval $[19/20,21/20]$ for the eigenvalues
of $K_E$. We can now compare the metric in \eqref{eq:9.10} with ordinary Frobenius
and Euclidean norms:
\begin{equation}\label{eq:9.11}
\|E\|_F^2+\frac{20}{21}\|j\|^2\le\frac{21}{40}\sigma.
\end{equation}

\textbf{Step 2: control the full perturbation.} Its squared norm contains
$2\|j\|^2$, so the estimate in \eqref{eq:9.11} needs one further bound on the coupling.
The Gram constraint gives $\|j\|^2\le x(t+1/40)\le47\sigma/160$.
The coefficient $47/160$ combines $xt\le9\sigma/32$ and
$x/40\le\sigma/80$. Also $x^2\le\sigma/4000$.
Combining these estimates leaves the convenient upper bound $5\sigma/6$:
\begin{equation}\label{eq:9.12}
\|\widehat F-I\|_F^2=x^2+2\|j\|^2+\|E\|_F^2\le\frac56\sigma,
\qquad \|\widehat F-I\|_{\rm op}<\frac1{30}.
\end{equation}

\textbf{Step 3: let the trace gain dominate the logarithmic remainder.}
The size assumption and \eqref{eq:9.12} put every eigenvalue perturbation in
$[-1/30,1/30]$. On this interval, integration of the derivative gives
$\log(1+z)\ge z-(15/29)z^2$; the coefficient is
$15/29=1/[2(1-1/30)]$.
The quadratic term is controlled by \eqref{eq:9.12}, while
$\operatorname{tr}(\widehat F-I)=(1-d/4)\sigma\ge\sigma/2$.
Their difference is
\begin{equation}\label{eq:9.13}
\log\det\widehat F\ge\frac\sigma2-\frac{15}{29}\frac{5\sigma}{6}
=\frac{2\sigma}{29}\ge\frac\sigma{16}.
\end{equation}

The bound $2\sigma/29$ is rounded down to $\sigma/16$ for subsequent use.
Finally, $\sigma=2\chi/(4-b^2)\ge\chi/2$. 
\end{proof}

\textbf{Completion of the deferred proof of Lemma~\ref{lem-geometric-control}.} The first two bounds in Lemma~\ref{lem-geometric-control}
are part of Lemma~\ref{lem-residual}. Its hypothesis is expressed in $\chi$, while the local
lemma uses $\sigma$, so we check the conversion explicitly. For the small-defect
assertion, $\chi\le3/2000$ implies
$\sigma=2\chi/(4-b^2)\le(8/15)\chi\le1/1250<10^{-3}$.
Lemma~\ref{lem-local} and \eqref{eq:8.2} therefore give $\Delta_{\rm geom}\ge\chi/32$.
This proves the remaining assertion and completes Lemma~\ref{lem-geometric-control}, as promised in Section~\ref{sec:6.1}.

\subsection{Remaining estimates for Theorem~\ref{thm-stability}: shape and center}\label{sec:9.4}

We now prove the polar-factor and support-mismatch estimates deferred in
Section~\ref{sec:6.2}. We finish the passage from certificate error to ellipsoid error under the
hypotheses of Theorem~\ref{thm-stability}, in particular $\epsilon\le10^{-6}$.
The contact estimate \eqref{eq:6.3} already gives $\chi\le32\epsilon$ and hence
$\|Q-M_0\|_F\le4\sqrt\epsilon$. What remains is to transfer this bound through
the polar factor and matrix inversion, then recover the center from support
values. The constants below are explicit bounds for these transfers.

\textbf{Step 1: keep the polar factor uniformly invertible.} Set
$\delta=\|Q-M_0\|_{\rm op}\le4\sqrt\epsilon\le0.004$. Since
$M_0=2vv^T+\sqrt{1-b}P$, the triangle inequality gives, for every vector $x$,
\begin{equation}\label{eq:9.14}
\left(\frac1{\sqrt2}-\delta\right)\|x\|
\le\|Q^Tx\|\le(2+\delta)\|x\|,
\qquad
\|Q^Tx\|^2=x^TM^2x.
\end{equation}

Apply these inequalities to unit eigenvectors of $M\succ0$. Every eigenvalue
lies in $[1/\sqrt2-\delta,\,2+\delta]$.
The rational bounds $7/10$ and $201/100$ lie outside this interval uniformly
under the assumed bound on $\delta$:
\begin{equation}\label{eq:9.15}
\lambda_{\min}(M)>\frac7{10},\qquad
\lambda_{\max}(M)<\frac{201}{100},\qquad
\|Q\|_{\rm op}\le2+\delta<\frac{201}{100}.
\end{equation}

The lower bound prevents matrix inversion from magnifying small errors
without control. The upper bound will control the product difference in the
next step.

\textbf{Step 2: transfer the error from $Q$ to its polar factor.} Since $M^2=QQ^T$,
the difference $X=M-M_0$ satisfies a linear matrix equation. Taking its
Frobenius inner product with $X$ uses the positive spectral lower bounds to
control $\|X\|_F$. The identity $MX+XM_0=QQ^T-M_0^2$ gives
\begin{equation}\label{eq:9.16}
\begin{aligned}
(\lambda_{\min}(M)+\lambda_{\min}(M_0))\|X\|_F^2
&\le\operatorname{tr}(X^TMX)+\operatorname{tr}(X^TXM_0)\\
&=\langle X,QQ^T-M_0^2\rangle_F\\
&\le\|X\|_F\|QQ^T-M_0^2\|_F.
\end{aligned}
\end{equation}

For $X=0$ the conclusion is immediate; otherwise cancel $\|X\|_F$. Since
$QQ^T-M_0^2=(Q-M_0)Q^T+M_0(Q^T-M_0)$,
\begin{equation}\label{eq:9.17}
\begin{aligned}
\|M-M_0\|_F
&\le\frac{\|QQ^T-M_0^2\|_F}
{\lambda_{\min}(M)+\lambda_{\min}(M_0)}\\
&\le\frac{(\|Q\|_{\rm op}+\|M_0\|_{\rm op})\|Q-M_0\|_F}{7/5}
\le\frac{401}{35}\sqrt\epsilon
\le12\sqrt\epsilon.
\end{aligned}
\end{equation}

Here $401/35$ is the coefficient obtained by inserting the spectral bounds;
$12$ is a convenient upper bound for it.

\textbf{Step 3: transfer the error to the ellipsoid shape.} Matrix matching
controls $A-M^{-1}$, and Step 2 controls $M^{-1}-M_0^{-1}$ using the same
positive lower bounds. The inverse-difference estimate in \eqref{eq:6.5} gives
\begin{equation}\label{eq:9.18}
\|A-A_*\|_F
\le\left(\frac{20}{7}+\frac{1200}{49}\right)\sqrt\epsilon
\le30\sqrt\epsilon.
\end{equation}

The two terms in the coefficient correspond to these two errors, and their
sum is less than $30$. The non-strict conclusion includes $\epsilon=0$.

\textbf{Step 4: express the model's support mismatch through the contact defect.}
To recover the center, we need to compare the model support value with the
original supporting plane at each contact. Squaring and subtracting exposes
the pole and latitude factors directly. For a unit vector $u=sv+y$, $y\perp v$,
compute
\begin{equation}\label{eq:9.19}
\|A_*u\|^2-\left(1+\frac s2\right)^2
=\frac{s^2}{4}+\frac{1-s^2}{1-b}-\left(1+\frac s2\right)^2
=\frac{(1+s)(b-s)}{1-b}.
\end{equation}

Rationalizing therefore gives
\begin{equation}\label{eq:9.20}
g_*(s):=\|A_*u\|-1-\frac s2
=\frac{(1+s)(b-s)}{(1-b)(\|A_*u\|+1+s/2)}.
\end{equation}

The squared mismatch is therefore the contact-defect summand times a bounded
factor. The bounds $\|A_*u\|+1+s/2\ge1$, $(1+s)(1-as)\le2$, and
$(1-b)^2\ge1/4$ give the coefficient $8$:
\begin{equation}\label{eq:9.21}
g_*(s)^2
\le8\,\frac{(1+s)(s-b)^2}{1-as},
\qquad
\sum_i\lambda_i g_*(s_i)^2\le8\chi.
\end{equation}

\textbf{Step 5: recover the center and collect the estimates.} Subtracting support
values expresses each projection of $c-c_*$ as the sum of the slack, the
model mismatch, and the shape error:
\begin{equation}\label{eq:9.22}
(c-c_*)^Tu_i=-r_i-g_*(s_i)+\|A_*u_i\|-\|Au_i\|.
\end{equation}

Multiplication by $\lambda_i u_i$ and summation uses isotropy to recover
$c-c_*$ itself. The triangle and Bessel estimates in Section~\ref{sec:6} give
\begin{equation}\label{eq:9.23}
\|c-c_*\|\le4\epsilon+\sqrt{8\chi}+\|A-A_*\|_F
\le50\sqrt\epsilon.
\end{equation}

In this last estimate, $\sqrt{8\chi}\le16\sqrt\epsilon$, the shape term is
at most $30\sqrt\epsilon$, and $4\epsilon\le0.004\sqrt\epsilon$.
The coefficient $50$ is a single convenient upper bound for their sum.
Together with the completed proof of Lemma~\ref{lem-geometric-control}, these estimates justify
all deferred inputs in the proof of Theorem~\ref{thm-stability}.
